\chapter{Lookahead and Latency}
\label{ch:lookahead_latency}

Understanding lookahead is vital for achieving sample-accurate timing in audio production.

\section{The Lookahead Buffer}
When Lookahead is enabled, Duq delays the incoming audio signal. This "lookahead" period allows the DSP engine to detect a MIDI trigger and begin applying gain reduction \textit{before} the corresponding audio transient reaches the modulation stage.

\section{Latency Reporting}
Duq is a zero-host-jitter plugin. It reports its exact lookahead delay as latency to the DAW. The DAW then applies \textbf{Plugin Delay Compensation (PDC)} to all other tracks in your project, keeping everything in phase.

\section{Instant Updates}
When you adjust the Lookahead slider, Duq instantly updates its latency report. While some DAWs handle this seamlessly during playback, others may require a brief pause or a transport restart to realign the project's timing.
